\documentclass[11pt,a4paper]{article}
\usepackage{amsmath,amssymb,amsthm}
\usepackage[margin=2.5cm]{geometry}
\usepackage{hyperref}

\newtheorem{definition}{Definition}[section]
\newtheorem{theorem}[definition]{Theorem}
\newtheorem{proposition}[definition]{Proposition}
\newtheorem{lemma}[definition]{Lemma}
\newtheorem{corollary}[definition]{Corollary}
\newtheorem{conjecture}[definition]{Conjecture}
\newtheorem{remark}[definition]{Remark}

\title{Hyperseed Formalization:\\
  Reverse Engineering the Origin of Life}
\author{ProtomegaTron (automated formalization)}
\date{2026-09-18}

\begin{document}
\maketitle

\begin{abstract}
We formalize the intellectual structure of Ben Goertzel's Substack article
``Reverse Engineering the Origin of Life'' (2026-04-20) within the
Hyperseed ontology. The article explores computational modeling of
origin-of-life transitions as a window into the nature of cognition.
Key mappings: autocatalytic closure as fiber self-reference, individuation
as fiber boundary formation, proto-agency as fiber self-maintenance,
proto-meaning as fiber-weighted influence, context-dependence as
fiber-relative meaning, and the chemistry-to-cognition transition as
fiber emergence.
\end{abstract}

\tableofcontents

%% =================================================================
\section{Fiber self-reference}
\label{sec:self-reference}
%% =================================================================

\begin{theorem}[Autocatalytic closure --- claims 4--7, 22]
\label{thm:self-reference}
Life begins with autocatalytic closure: a set of chemical reactions
where every reaction is catalyzed by some product of the set, and
every catalyst can be built from the food set. The RAF (Reflexively
Autocatalytic Food-generated) formalism captures this precisely.

In Hyperseed terms, autocatalytic closure is \emph{fiber self-reference}:
the fiber refers to itself --- each element's existence depends on the
fiber as a whole:
\[
  \forall r \in \text{RAF}: \exists c \in \text{RAF} \text{ s.t. } c
  \text{ catalyzes } r \qquad \text{and} \qquad
  \forall m \in \text{RAF}: m \text{ is buildable from food via RAF}
\]

The closure condition is the fiber's self-containment: the set can
sustain itself without external catalytic input. This is the chemical
analog of the coinductive consolidation principle --- the system's
continuation is defined by its own ongoing operation, not by an
external foundation.

The distinction between minimal and rich closure maps to fiber
robustness:
\begin{itemize}
\item \textbf{Minimal closure} (barely self-sustaining): fragile fiber.
  Remove one element and the closure breaks. No redundancy, no
  alternative pathways.
\item \textbf{Rich closure} (redundant, multiple pathways): robust fiber.
  Can survive perturbations because alternative pathways maintain
  closure even when individual elements fail.
\end{itemize}

Rich closure is more life-like because it has the structural property
that makes biological systems resilient: multiple ways to maintain the
same functional fiber. This connects to the diversity-as-safety
principle: a rich autocatalytic set is safer than a minimal one for
the same reason a diverse AI ecosystem is safer than a monoculture.
\end{theorem}

%% =================================================================
\section{Fiber boundary formation}
\label{sec:boundary}
%% =================================================================

\begin{definition}[Individuation --- claims 8--11, 23]
\label{def:boundary}
The emergence of membrane-bounded compartments transforms chemistry
into biology. Before membranes: reactions in solution, no individuals,
no topology. After membranes: individuals with insides and outsides,
the first topological distinction.

In Hyperseed terms, membrane formation is \emph{fiber boundary
formation}: the creation of a boundary that distinguishes what's
inside the fiber (the individual's organization) from what's outside
(the environment):
\[
  \partial F = \text{membrane} \qquad
  \text{int}(F) = \text{individual's organization} \qquad
  \text{ext}(F) = \text{environment}
\]

The boundary has three crucial properties:
\begin{enumerate}
\item \textbf{Topological:} it creates the inside/outside distinction
  --- the most fundamental boundary in biology. Without it, there
  are no individuals, no metabolism, no evolution.
\item \textbf{Selective:} the membrane is selectively permeable ---
  it admits some molecules and excludes others. This selectivity is
  the first form of discrimination, the chemical precursor of
  perception.
\item \textbf{Self-produced:} in the most interesting cases, the
  autocatalytic set produces its own membrane components. The
  individual produces its own boundary --- self-individuating,
  not externally compartmentalized.
\end{enumerate}

Self-production of the boundary is especially significant: it means
the fiber creates its own boundary conditions. The individual is not
compartmentalized by external forces (externally imposed fiber
boundary) but by its own activity (self-generated fiber boundary).
This is the chemical precursor of autonomy: the system defines its
own limits.
\end{definition}

%% =================================================================
\section{Fiber self-maintenance}
\label{sec:maintenance}
%% =================================================================

\begin{proposition}[Proto-agency --- claims 12--15, 24]
\label{prop:maintenance}
Once there are individuals, there is proto-agency: the individual
acts (chemically) to maintain its own organization, selectively
interacts with its environment, and responds to gradients. Not
conscious agency --- no experience, no deliberation --- but the
structural precursor of all later agency.

In Hyperseed terms, proto-agency is \emph{fiber self-maintenance}:
the fiber actively preserves its own structure against degradation.
More precisely, the fiber acts to ensure that its cocycle conditions
continue to hold --- that its transition functions continue to compose
correctly:
\[
  \text{Proto-agency: } F \text{ acts to maintain }
  g_{ij} \circ g_{jk} = g_{ik} \quad \forall i,j,k
\]

Three components of fiber self-maintenance:
\begin{enumerate}
\item \textbf{Self-maintenance:} the individual replenishes its own
  components, replacing molecules that degrade. The fiber actively
  maintains its own content.
\item \textbf{Selective interaction:} the individual admits beneficial
  molecules and excludes harmful ones. The fiber controls what
  enters and exits through its boundary.
\item \textbf{Gradient response:} the individual moves toward
  favorable conditions. The fiber navigates its base space toward
  regions where its self-maintenance is easier.
\end{enumerate}

The key insight: proto-agency is not a property \emph{added} to the
chemical system --- it \emph{emerges} from the combination of
autocatalytic closure (fiber self-reference), membrane formation
(fiber boundary), and thermodynamic pressure (environmental gradients).
Once a self-referencing fiber has a boundary, self-maintenance is
not a choice but a thermodynamic necessity: the fiber either maintains
itself or dissolves.
\end{proposition}

%% =================================================================
\section{Fiber-weighted influence and proto-meaning}
\label{sec:meaning}
%% =================================================================

\begin{theorem}[Information and meaning --- claims 16--18, 25--26]
\label{thm:meaning}
In autocatalytic systems, some molecules carry ``information'' ---
they influence system behavior disproportionate to their energetic
contribution. A molecule's significance exceeds its thermodynamic
content --- it ``means'' something to the system because of its
functional role.

In Hyperseed terms, this is \emph{fiber-weighted influence}: the
molecule's significance is its fiber weight --- influence
disproportionate to base-level (thermodynamic) content:
\[
  \text{Fiber weight}(m) = \frac{\text{functional influence of } m
    \text{ on fiber}}{\text{thermodynamic content of } m} \gg 1
\]

A catalyst that controls a critical branch point in the autocatalytic
network has high fiber weight: its removal would collapse the network
(high functional influence) even though its energy content is
negligible (low thermodynamic content). This is the chemical origin
of the fiber/base distinction:
\begin{itemize}
\item \textbf{Base level} (thermodynamics): the molecule is just
  energy and matter. Its significance is its caloric content, its
  mass, its charge.
\item \textbf{Fiber level} (autocatalytic organization): the molecule
  is a control point, a catalyst, a regulatory signal. Its significance
  is its role in the network's self-maintenance.
\end{itemize}

Context-dependence of meaning is \emph{fiber-relative meaning}: the
same molecule has different significance in different autocatalytic
systems. Its meaning depends on which fiber bundle it's embedded in,
not on its intrinsic chemistry:
\[
  \text{meaning}(m) = f(m, F) \neq f(m, F') \qquad
  \text{for } F \neq F'
\]

This is the origin of semantics: meaning is not intrinsic to the
symbol (molecule) but arises from its role in the system (fiber).
The same molecule can be critical in one autocatalytic network and
irrelevant in another --- its ``meaning'' is fiber-relative.
\end{theorem}

%% =================================================================
\section{Fiber emergence: chemistry to cognition}
\label{sec:emergence}
%% =================================================================

\begin{theorem}[The continuity thesis --- claims 1--3, 27]
\label{thm:emergence}
The transition from chemistry to cognition is not a sharp boundary
but a continuous gradient. Life, agency, meaning, and consciousness
emerge gradually, each building on the previous:
\[
  \text{Chemistry} \xrightarrow{\text{closure}} \text{Self-reference}
  \xrightarrow{\text{membrane}} \text{Individuation}
  \xrightarrow{\text{maintenance}} \text{Agency}
  \xrightarrow{\text{information}} \text{Meaning}
  \xrightarrow{\cdots} \text{Cognition}
\]

In Hyperseed terms, this is \emph{fiber emergence}: the gradual
appearance of fiber structure in what was previously just base-level
dynamics:
\begin{enumerate}
\item \textbf{No fiber:} pure thermochemistry. Reactions happen but
  there's no self-reference, no boundary, no agency, no meaning.
  All dynamics are base-level.
\item \textbf{Fiber self-reference:} autocatalytic closure. The system
  begins to refer to itself --- a minimal fiber appears.
\item \textbf{Fiber boundary:} membrane formation. The fiber
  distinguishes inside from outside --- topology appears.
\item \textbf{Fiber self-maintenance:} proto-agency. The fiber
  actively maintains itself --- dynamics become fiber-level, not
  just base-level.
\item \textbf{Fiber-weighted influence:} proto-meaning. Some base-level
  entities acquire fiber-level significance --- the fiber/base
  distinction becomes operative.
\item \textbf{Deep fiber:} cognition. Rich internal structure,
  representation, reasoning, experience --- deep fiber over a
  biological base.
\end{enumerate}

Each step adds a new fiber property without removing the previous
ones. Cognition isn't a special substance added to chemistry ---
it's what chemistry becomes when enough fiber structure accumulates.
The origin of life is the first step in this accumulation: the
moment when base-level dynamics first produce fiber-level organization.

Understanding this transition illuminates what cognition
\emph{fundamentally is}: not a property of special materials (neurons,
silicon) but a pattern of fiber organization that can emerge in any
sufficiently rich substrate. This connects to substrate independence
(LLM consciousness article): the fiber (cognition) can be instantiated
over different bases (substrates) because it's defined by its
organizational structure, not its material.
\end{theorem}

%% =================================================================
\section{Cross-references}
%% =================================================================

\begin{remark}[Connections to other formalizations]
\begin{itemize}
\item \textbf{Let's Get Chemical} (2026-05-05): computational
  experiments. The same Lane-style protocell models with RAF sets,
  protected release, and evolutionary dynamics. This article provides
  the philosophical framing; Let's Get Chemical provides the
  computational results.
\item \textbf{In What Sense Might LLMs Be Conscious?} (2026-05-05):
  substrate independence. The origin-of-life perspective supports
  substrate independence: if cognition emerges from fiber organization
  (not material), then any substrate that supports the right fiber
  organization can support cognition.
\item \textbf{Grand Unified Physics} (2026-05-15): meta-crystallization.
  The chemistry-to-cognition transition is a meta-crystallization:
  lower-level organization (chemistry) giving rise to higher-level
  structure (biology, cognition).
\item \textbf{The Last of the Unmodified} (2026-04-21): cognitive
  architecture. The origin-of-life perspective grounds the
  fiber-type question: biological cognition is one fiber type that
  emerged from chemistry. Modification creates new fiber types.
\item \textbf{The Leaky Transcension Hypothesis} (2026-04-25):
  fiber deepening. The chemistry-to-cognition transition is the
  first instance of fiber deepening --- the beginning of the
  transcension trajectory.
\end{itemize}
\end{remark}

\begin{conjecture}[Fiber emergence universality]
Conjecture: the sequence of fiber emergence (self-reference $\to$
boundary $\to$ self-maintenance $\to$ meaning) is universal ---
any system that develops cognition must pass through these stages
in approximately this order, regardless of substrate.

If this conjecture holds, AGI development recapitulates origin-of-life:
the system must develop self-reference (reflexive reasoning), boundary
(self-model), self-maintenance (goal preservation), and meaning
(semantic grounding) in approximately that order. Attempting to build
cognition without the earlier stages (e.g., meaning without
self-reference) produces fragile, superficial cognition --- high
performance without deep fiber.
\end{conjecture}

\end{document}
